\section{Makroökonomisches Modell}

\subsection{Grundidee}

Das makroökonomische Modell beschreibt die gesamtwirtschaftliche
Gleichgewichtssituation einer Volkswirtschaft.

Das Gleichgewicht entsteht durch das Zusammenspiel von

\begin{itemize}
    \item aggregierter Nachfrage (AN)
    \item aggregiertem Angebot (AA)
\end{itemize}

Das Modell bestimmt zwei zentrale Grössen:

\begin{itemize}
    \item Preisniveau
    \item reales Bruttoinlandprodukt (BIP)
\end{itemize}


\subsection{Achsen des Modells}

\textbf{Horizontale Achse}

Reales Bruttoinlandprodukt (BIP)

\[
    BIP_{real} = \text{inflationsbereinigte Wertschöpfung eines Landes}
\]

Definition Wertschöpfung:

\[
    \text{Wertschöpfung} =
    \text{Wert aller produzierten Güter}
    -
    \text{Vorleistungen}
\]

Messgrösse: Geldbetrag.

\textbf{Vertikale Achse}

Preisniveau

Messung über Preisindizes, z.B.

\begin{itemize}
\item Landesindex der Konsumentenpreise (LIK)
\item Produzentenpreisindex
\end{itemize}

Preisniveau wird meist als Index dargestellt:

\[
    Preisniveau_{Basisjahr} = 100
\]


\subsection{Aggregierte Nachfrage (AN)}

Die aggregierte Nachfrage beschreibt die gesamte Güternachfrage einer Volkswirtschaft.

Sie setzt sich aus vier Nachfragekomponenten zusammen:

\[
    AN = C + I + G + NX
\]

mit

\begin{itemize}
    \item $C$ = Konsum der Haushalte
    \item $I$ = Investitionen der Unternehmen
    \item $G$ = Staatsausgaben
    \item $NX$ = Nettoexporte (Exporte – Importe)
\end{itemize}

Eigenschaft:

\[
    \text{Je tiefer das Preisniveau, desto höher die Nachfrage}
\]

Daher ist die AN-Kurve fallend.


\subsection{Kurzfristige aggregierte Angebotskurve (AAK)}

Die kurzfristige Angebotskurve zeigt das Güterangebot der Unternehmen.

Kurzfristig sind Produktionsfaktoren gegeben:

\begin{itemize}
    \item Arbeit
    \item Kapital
    \item Technologie
    \item natürliche Ressourcen
\end{itemize}

Eigenschaft:

\[
    \text{Je höher das Preisniveau, desto höher das Angebot}
\]

Die AAK ist deshalb steigend.


\subsection{Langfristige Angebotskurve (AAL)}

Die langfristige Angebotskurve beschreibt die maximale nachhaltige
Produktionskapazität einer Volkswirtschaft.

Diese Grenze nennt man

\[
    \textbf{Kapazitätsgrenze}
\]

An diesem Punkt sind die Produktionsfaktoren optimal ausgelastet.

Die AAL ist daher \textbf{vertikal}.


\subsection{Makroökonomisches Gleichgewicht}

Das gesamtwirtschaftliche Gleichgewicht entsteht beim Schnittpunkt von

\[
    AN = AAK = AAL
\]

In diesem Punkt gilt:

\begin{itemize}
    \item Produktionsfaktoren sind optimal ausgelastet
    \item reales BIP entspricht dem Produktionspotential
    \item Preisniveau ist stabil
\end{itemize}


\subsection{Über- und Unterauslastung}

\textbf{Überauslastung}

\begin{itemize}
    \item Produktion über Kapazitätsgrenze
    \item z.B. durch Überstunden oder maximale Auslastung
    \item führt häufig zu Inflation
\end{itemize}

\textbf{Unterauslastung}

\begin{itemize}
    \item Produktion unter Kapazitätsgrenze
    \item ungenutzte Produktionsfaktoren
    \item führt zu konjunktureller Arbeitslosigkeit
\end{itemize}